\section{Multiplication: convolution of logarithmic blocks}
\label{plt:sec:multiplication}

Multiplication is the operation on which every later construction rests:
reciprocals, powers, composition, and reversion are all built from it by
triangular recurrences.  Its whole content is that the Cauchy rule in the
outer variable~$t$ is \emph{coefficientwise finite}, so that each output
block is an exact finite expression in finitely many input blocks.  Nothing
here is an analytic rearrangement, and nothing here requires convergence.

Throughout this section $\K$ is a field of characteristic zero, $X$ is the
large variable, the regime is $X\to+\infty$ on the positive ray unless a
different domain is announced, and
\begin{equation}\label{plt:eq:mul-generators}
 t\DefinitionEquals X^{-1},
 \qquad
 L\DefinitionEquals\log X .
\end{equation}
In the formal algebra $t$ and $L$ are independent generators; their analytic
coupling is recorded only by the distinguished derivation, which plays no
role in this section.  The four ambient classes are

\begin{center}
\begin{tabular}{@{}lll@{}}
\hline
class & coefficient ring & elements \\
\hline
$\PLpow$ & $\K[L]$ & $\sum_{n\ge0}p_n(L)t^n$ \\
$\PLlau$ & $\K[L]$ & $\sum_{n\ge n_0}p_n(L)t^n$, $n_0\in\IntegerNumbers$ \\
$\PLrat$ & $\K(L)$ & $\sum_{n\ge n_0}p_n(L)t^n$, $p_n$ rational in $L$ \\
$\PLit$  & $\K((L^{-1}))$ & $\sum_{n\ge n_0}t^n\sum_{j\ge j_n}a_{n,j}L^{-j}$ \\
\hline
\end{tabular}
\end{center}

All truncations below are \emph{exclusive}: $\Trunc{F}{N}$ keeps every
logarithmic monomial in every block $t^n$ with $n<N$, and discards the rest.

\begin{remark}[Reconciliation of the source conventions]
\label{plt:rmk:mul-conventions}
Two of the merged reports use an inclusive outer cutoff $[F]_{\le N}$.  The
translation is $[F]_{\le N}=\Trunc{F}{N+1}$, and every dependency range,
algorithm bound and error exponent quoted from those reports has been shifted
by one block accordingly.  Likewise, one report writes the coefficient
bracket as $[F]_M$; here the series always stands to the right of the
bracket, $[t^n]F=p_n(L)$ and $[t^nL^j]F=a_{n,j}$.
\end{remark}

\subsection{The block Cauchy product}
\label{plt:sec:mul-cauchy}

It is economical to prove the algebra of all four classes at once.  What the
three coefficient rings $\K[L]$, $\K(L)$, $\K((L^{-1}))$ have in common is a
degree function in $L$ together with a coefficient-extraction map obeying the
usual convolution rule.

\begin{definition}[Logarithmically graded domain]
\label{plt:def:mul-graded-domain}
A \emph{logarithmically graded domain over $\K$} is a commutative
$\K$-algebra $R$ together with $\K$-linear maps
$\kappa_d\colon R\to\K$, one for each $d\in\IntegerNumbers$, such that
for every nonzero $f\in R$ the set $\SetOf{d:\kappa_d(f)\ne0}$ is nonempty
and bounded above, and:
\begin{enumerate}
\item[(G1)] $f=0$ if and only if $\kappa_d(f)=0$ for every $d$;
\item[(G2)] $\kappa_d(fg)=\sum_{i+j=d}\kappa_i(f)\kappa_j(g)$ for all
$f,g\in R$ and all $d\in\IntegerNumbers$, the sum having only finitely many
nonzero terms;
\item[(G3)] $\kappa_0(1)=1$ and $\kappa_d(1)=0$ for $d\ne0$.
\end{enumerate}
For $f\ne0$ put $\degL f\DefinitionEquals\max\SetOf{d:\kappa_d(f)\ne0}$ and
$\lc_L(f)\DefinitionEquals\kappa_{\degL f}(f)\in\K^\times$; set
$\degL 0\DefinitionEquals-\infty$ and $\lc_L(0)\DefinitionEquals0$.
\end{definition}

Here $\lc_L$ is the leading coefficient \emph{inside one block}; it is not
the same object as the leading scalar $\lc$ of a series, defined in
\cref{plt:def:mul-leading-data} below, although the two agree on the leading
block.

\begin{lemma}[Elementary consequences]
\label{plt:lem:mul-graded-basics}
Let $R$ be a logarithmically graded domain over $\K$.  For all $f,g\in R$:
\begin{equation}\label{plt:eq:mul-degL-laws}
 \degL(fg)=\degL f+\degL g,
 \qquad
 \lc_L(fg)=\lc_L(f)\lc_L(g),
 \qquad
 \degL(f+g)\le\max(\degL f,\degL g).
\end{equation}
Here $-\infty+c=-\infty$ and $\max(-\infty,c)=c$.  In particular $R$ is an
integral domain.  Moreover, if $\degL f\le d$ and $\degL g\le e$, then
$\kappa_{d+e}(fg)=\kappa_d(f)\kappa_e(g)$.
\end{lemma}

\begin{proof}
If $f=0$ or $g=0$ every assertion is trivial under the stated conventions.
The third assertion in \eqref{plt:eq:mul-degL-laws} is immediate from
$\K$-linearity of each $\kappa_d$.  For the first two, assume $f,g\ne0$ and
set $d=\degL f$, $e=\degL g$.  If $c>d+e$ and $i+j=c$, then $i>d$ or $j>e$, so
every term of (G2) vanishes and $\kappa_c(fg)=0$.  If $c=d+e$, the only
possibly nonvanishing term is $i=d$, $j=e$, whence
\[
 \kappa_{d+e}(fg)=\kappa_d(f)\kappa_e(g)=\lc_L(f)\lc_L(g)\ne0,
\]
because $\K$ is a field and both factors are nonzero.  This proves
$\degL(fg)=d+e$, $\lc_L(fg)=\lc_L(f)\lc_L(g)$, and, by (G1), $fg\ne0$; so $R$
is a domain.  The last sentence is the same computation with $d,e$ merely
upper bounds, all other terms of (G2) again vanishing.
\end{proof}

\begin{lemma}[The three coefficient rings]
\label{plt:lem:mul-three-rings}
Each of
\[
 \K[L],\qquad \K(L),\qquad \K((L^{-1}))
\]
is a logarithmically graded domain over $\K$, with $\degL$ equal
respectively to the polynomial degree, to $\degL(p/q)=\deg p-\deg q$, and to
the largest exponent occurring in the Laurent expansion.  There are injective
$\K$-algebra homomorphisms
\begin{equation}\label{plt:eq:mul-coeff-chain}
 \K[L]\hookrightarrow\K(L)\hookrightarrow\K((L^{-1}))
\end{equation}
that preserve $\degL$ and $\lc_L$.
\end{lemma}

\begin{proof}
For $\K[L]$ take $\kappa_d(f)$ to be the coefficient of $L^d$ in $f$ for
$d\ge0$ and $\kappa_d=0$ for $d<0$; (G1)--(G3) are the definition of
polynomial multiplication.  For $\K((L^{-1}))$, an element is
$f=\sum_{j\le d}a_jL^j$ with only finitely many $j>0$ allowed but arbitrarily
negative $j$ permitted; put $\kappa_j(f)=a_j$.  Then (G1) and (G3) are clear,
and in (G2) the constraint $i+j=c$ with $i\le\degL f$, $j\le\degL g$ forces
$i\ge c-\degL g$, so the sum is finite; it is the definition of the Laurent
product.  The support of $f$ is bounded above by construction.

The inclusion $\K[L]\subset\K((L^{-1}))$ is a $\K$-algebra map, and every
nonzero element of $\K((L^{-1}))$ is invertible: writing $f\ne0$ as
$f=a_dL^d(1+u)$ with $d=\degL f$ and $u\in L^{-1}\K[[L^{-1}]]$, the family
$\bigl((-u)^r\bigr)_{r\ge0}$ is $L^{-1}$-adically summable and its sum
inverts $1+u$, so $f^{-1}=a_d^{-1}L^{-d}\sum_{r\ge0}(-u)^r$.  Hence
every nonzero polynomial becomes invertible, and the universal property of the
field of fractions supplies a unique $\K$-algebra homomorphism
$\iota\colon\K(L)\to\K((L^{-1}))$ extending the inclusion; it is injective
because $\K(L)$ is a field and $\iota(1)=1$.  Pulling the maps $\kappa_d$ back
along $\iota$ makes $\K(L)$ a logarithmically graded domain, and
\cref{plt:lem:mul-graded-basics} applied to $\iota(p)=\iota(q)\iota(p/q)$
gives $\degL(p/q)=\deg p-\deg q$ and
$\lc_L(p/q)=\lc_L(p)\lc_L(q)^{-1}$.  Both maps in
\eqref{plt:eq:mul-coeff-chain} preserve $\kappa_d$ by construction, hence
preserve $\degL$ and $\lc_L$.
\end{proof}

\begin{definition}[The block Cauchy product]
\label{plt:def:mul-cauchy}
Let $R$ be a logarithmically graded domain over $\K$ and let
\begin{equation}\label{plt:eq:mul-FG-blocks}
 F=\sum_{n\ge a}p_n\,t^n,
 \qquad
 G=\sum_{m\ge b}q_m\,t^m,
 \qquad p_n,q_m\in R,
\end{equation}
be elements of $R((t))$, the ring of formal Laurent series in $t$ over $R$.
Their product is
\begin{equation}\label{plt:eq:mul-cauchy}
 FG\DefinitionEquals\sum_{N\ge a+b}C_N\,t^N,
 \qquad
 C_N\DefinitionEquals\sum_{\substack{n+m=N\\ n\ge a,\ m\ge b}}p_nq_m
 =\sum_{n=a}^{N-b}p_n\,q_{N-n}.
\end{equation}
\end{definition}

\begin{lemma}[Finiteness of every fiber]
\label{plt:lem:mul-finite-fiber}
For fixed $N$ the index set in \eqref{plt:eq:mul-cauchy} is
$\SetOf{n\in\IntegerNumbers:a\le n\le N-b}$, of cardinality
$\max(N-a-b+1,\,0)$.  In particular it is finite, and it is empty exactly
when $N<a+b$.
\end{lemma}

\begin{proof}
The constraints $n\ge a$ and $m=N-n\ge b$ are equivalent to
$a\le n\le N-b$.  An interval of integers $[a,N-b]$ has $N-b-a+1$ elements
when $N-b\ge a$ and is empty otherwise.
\end{proof}

This is the entire mechanism.  There is no limit, no rearrangement, and no
hypothesis of numerical convergence: each $C_N$ is a finite sum of products
in~$R$.

\begin{theorem}[Closure of the four ambient classes]
\label{plt:thm:mul-closure}
Let $R$ be a logarithmically graded domain over $\K$.  Then
\eqref{plt:eq:mul-cauchy} makes $R((t))$ a commutative $\K$-algebra with unit
$1$, containing $R[[t]]$ as a subalgebra.  Consequently each of
\[
 \PLpow=\K[L][[t]],\quad
 \PLlau=\K[L]((t)),\quad
 \PLrat=\K(L)((t)),\quad
 \PLit=\K((L^{-1}))((t))
\]
is closed under multiplication, and the chain of inclusions
\begin{equation}\label{plt:eq:mul-class-chain}
 \PLpow\subset\PLlau\subset\PLrat\subset\PLit
\end{equation}
consists of injective $\K$-algebra homomorphisms.
\end{theorem}

\begin{proof}
By \cref{plt:lem:mul-finite-fiber} each $C_N$ is a well-defined element
of~$R$, and $C_N=0$ for $N<a+b$; so $FG$ again has support bounded below and
lies in $R((t))$.  Commutativity is the symmetry of the index set under
$(n,m)\mapsto(m,n)$; bilinearity over $\K$ is clear; and $1\cdot F=F$ because
the only admissible pair for $1=1\cdot t^0$ at level $N$ is $(0,N)$.  For
associativity let $H=\sum_{k\ge c}h_kt^k$.  Then
\[
 [t^N]\bigl((FG)H\bigr)
 =\sum_{r+k=N}\Bigl(\sum_{n+m=r}p_nq_m\Bigr)h_k
 =\sum_{n+m+k=N}p_nq_mh_k,
\]
the last sum ranging over the finite set of triples with $n\ge a$, $m\ge b$,
$k\ge c$: indeed $n+m+k=N$ with the three lower bounds confines each index to
a finite interval, so the regrouping is a finite rearrangement.  The same
computation with the other bracketing gives the same finite sum, whence
$(FG)H=F(GH)$.  Distributivity is likewise a finite regrouping.

If $\ordt F\ge0$ and $\ordt G\ge0$, then $C_N=0$ for $N<0$, so $R[[t]]$ is
closed under the product; it is a subalgebra.  Taking $R=\K[L]$ gives the
first two classes, $R=\K(L)$ the third, $R=\K((L^{-1}))$ the fourth, using
\cref{plt:lem:mul-three-rings}.  Finally, a $\K$-algebra homomorphism
$\varphi\colon R\to R'$ induces a $\K$-algebra homomorphism
$R((t))\to R'((t))$ by $\sum p_nt^n\mapsto\sum\varphi(p_n)t^n$, because
$\varphi$ commutes with the finite sums \eqref{plt:eq:mul-cauchy}; injectivity
is inherited blockwise.  Applying this to
\eqref{plt:eq:mul-coeff-chain} yields \eqref{plt:eq:mul-class-chain}.
\end{proof}

\begin{proposition}[The monomial-level convolution and the support bound]
\label{plt:prop:mul-monomial}
Let $F=\sum_{n,j}a_{n,j}t^nL^j$ and $G=\sum_{m,k}b_{m,k}t^mL^k$ lie in
$\PLlau$ (so $j,k\in\NaturalNumbers$, and $\PLpow$ is the case
$\ordt\ge0$) or in $\PLit$ (so $j,k\in\IntegerNumbers$, bounded above at
each fixed outer level).  Then
\begin{equation}\label{plt:eq:mul-2d-convolution}
 [t^NL^d](FG)
 =\sum_{n+m=N}\ \sum_{i+j=d}a_{n,i}b_{m,j},
\end{equation}
a finite sum, and
\begin{equation}\label{plt:eq:mul-support}
 \CoefficientSupportOf{FG}
 \subseteq
 \CoefficientSupportOf{F}+\CoefficientSupportOf{G},
\end{equation}
the Minkowski sum in $\IntegerNumbers^2$.
\end{proposition}

\begin{proof}
Formula \eqref{plt:eq:mul-2d-convolution} is (G2) applied to
\eqref{plt:eq:mul-cauchy}.  Finiteness holds because the outer sum is finite
by \cref{plt:lem:mul-finite-fiber} and, for each of its terms, the inner sum
is finite by (G2).  For \eqref{plt:eq:mul-support}: if
$(N,d)\notin\CoefficientSupportOf{F}+\CoefficientSupportOf{G}$, then every
summand of \eqref{plt:eq:mul-2d-convolution} has $a_{n,i}=0$ or $b_{m,j}=0$,
so $[t^NL^d](FG)=0$.
\end{proof}

\begin{remark}[$\PLrat$ has no monomial support of its own]
\label{plt:rmk:mul-rat-no-support}
\cref{plt:prop:mul-monomial} is deliberately not stated for $\PLrat$.  An
element of $\K(L)((t))$ has coefficient blocks that are rational functions of
$L$, and a rational function is not a (finite or infinite) combination of
powers of $L$ until one chooses an expansion point.  The coefficient support
$\CoefficientSupportOf{\cdot}$ becomes defined only after applying the
embedding $\PLrat\hookrightarrow\PLit$ of \cref{plt:thm:mul-closure}, which
expands each block at $L=\infty$ and turns the exact coefficient $1/(L+1)$
into the infinite series $L^{-1}-L^{-2}+L^{-3}-\cdots$.  Exact rational
dependence on $L$ and a monomial support in $L$ are different pieces of data;
only the second is what \eqref{plt:eq:mul-support} constrains.
\end{remark}

\begin{remark}[Why multiplication is a two-level, not a flat, operation]
\label{plt:rmk:mul-two-level}
Although \eqref{plt:eq:mul-2d-convolution} is a genuine two-dimensional
convolution, it should not be implemented as one.  The outer index is the
asymptotic hierarchy and the inner index is not: at fixed $N$ the inner sum
is bounded, whereas a cutoff imposed jointly on $|N|+|d|$ does not respect
dominance.  See \cref{plt:rmk:mul-rectangular-cutoff}.
\end{remark}

\begin{remark}[Hahn supports]
\label{plt:rmk:mul-hahn}
Nothing in \cref{plt:lem:mul-finite-fiber} uses that the exponents are
integers, only that the two supports are well ordered in the direction of
increasing smallness.  If $S,T\subseteq\Gamma$ are reverse-well-ordered
subsets of an ordered abelian group, then $S+T$ is reverse-well-ordered and
each $\gamma\in S+T$ has only finitely many decompositions $\gamma=s+t$ with
$s\in S$, $t\in T$; this is the Hahn support theorem
\cite{hahn,edgar-beginners}.  Everything in this section that does not
mention the specific value group $\IntegerNumbers$ transfers verbatim to the
Hahn power-log extension $\sum_{\alpha\in S}p_\alpha(L)t^\alpha$ and, in
particular, to the Puiseux-log extensions $\K[L]((t^{1/s}))$ used later for
fractional powers.  We do not repeat the statements.
\end{remark}

\subsection{Finite determination and precision bookkeeping}
\label{plt:sec:mul-finite-det}

The content of \cref{plt:lem:mul-finite-fiber} deserves to be recorded as a
dependency statement, because that is the form in which it is used: it says
exactly how much of the input must be known in order to know a given part of
the output, and it says that no more is needed and no less will do.

\begin{theorem}[Finite determination, with exact dependency range]
\label{plt:thm:mul-finite-determination}
Let $R$ be a logarithmically graded domain over $\K$, let $a,b,N\in
\IntegerNumbers$, and let $F,G\in R((t))$ satisfy $\ordt F\ge a$ and
$\ordt G\ge b$.  Then
\begin{equation}\label{plt:eq:mul-determination}
 [t^N](FG)=\sum_{n=a}^{N-b}p_n\,q_{N-n},
\end{equation}
so that $[t^N](FG)$ is a $\K$-bilinear function of the two finite tuples
\[
 (p_a,p_{a+1},\ldots,p_{N-b})
 \qquad\text{and}\qquad
 (q_b,q_{b+1},\ldots,q_{N-a}),
\]
equivalently of $\Trunc{F}{N-b+1}$ and $\Trunc{G}{N-a+1}$.  Both ranges are
exact: for every $n_0$ with $a\le n_0\le N-b$ there are $F,F'\in R((t))$ with
$\ordt\ge a$ that differ only in the block of index $n_0$, and a $G$ with
$\ordt G\ge b$, such that $[t^N](FG)\ne[t^N](F'G)$; and symmetrically in the
second argument.
\end{theorem}

\begin{proof}
Equation \eqref{plt:eq:mul-determination} and the bilinearity are
\cref{plt:def:mul-cauchy} together with \cref{plt:lem:mul-finite-fiber}; the
blocks listed are precisely those occurring in the sum.  For exactness fix
$n_0\in[a,N-b]$ and put $m_0=N-n_0$, so that $m_0\ge b$.  Take
$G=t^{m_0}$, which has $\ordt G=m_0\ge b$.  If $n_0>a$ set $F=t^a+t^{n_0}$
and $F'=t^a$; if $n_0=a$ set $F=t^a$ and $F'=2t^a$ (here $2\ne0$ because
$\operatorname{char}\K=0$).  In both cases $\ordt F=\ordt F'=a$, the two
series agree outside index $n_0$, and $[t^N](FG)-[t^N](F'G)$ equals $1$
respectively $-1$.  The symmetric statement follows by commutativity.
\end{proof}

\begin{checkpointbox}
\textbf{The precision rule.}  If $F$ is known through $\Trunc{F}{P}$ and $G$
through $\Trunc{G}{Q}$, and $\ordt F=a$, $\ordt G=b$, then $FG$ is known
through
\[
 \Trunc{FG}{\min(P+b,\;Q+a)}
\]
and, in the worst case, no further.  The binding entry is contributed by
whichever factor is paired with the \emph{more singular} partner: a factor
known to $P$ blocks is degraded by the valuation $b$ of the other factor.  In
particular, multiplying by a series with a pole of order $r$, that is with
valuation $-r$, costs exactly $r$ blocks of precision.
\end{checkpointbox}

\begin{corollary}[Truncated product]
\label{plt:cor:mul-truncated-product}
Let $F,G\in R((t))$ be nonzero with $a=\ordt F$, $b=\ordt G$, and let
$N\in\IntegerNumbers$.  Then for all integers $P\ge N-b$ and $Q\ge N-a$,
\begin{equation}\label{plt:eq:mul-truncated-product}
 \Trunc{FG}{N}
 =\Trunc{\bigl(\Trunc{F}{P}\bigr)\bigl(\Trunc{G}{Q}\bigr)}{N};
\end{equation}
% ed.: the optimality of $(N-b,N-a)$ was asserted unconditionally.  For
% ed.: $N\le a+b$ both sides of \eqref{plt:eq:mul-truncated-product} vanish
% ed.: whatever $P,Q$ are, so no pair is forced and there is nothing to be
% ed.: optimal about; the clause is now stated under $N>a+b$, which is what
% ed.: the argument actually establishes.
and if $N>a+b$ the pair $(P,Q)=(N-b,\,N-a)$ is optimal: neither entry can be
lowered by one, that is, there exist $F,G$ with these valuations
for which $P=N-b-1$, respectively $Q=N-a-1$, falsifies
\eqref{plt:eq:mul-truncated-product}.  If $F,G\in\PLpow$ and $N\ge0$, then
$(P,Q)=(N,N)$ is admissible, so the cruder law
\begin{equation}\label{plt:eq:mul-truncated-product-crude}
 \Trunc{FG}{N}=\Trunc{\bigl(\Trunc{F}{N}\bigr)\bigl(\Trunc{G}{N}\bigr)}{N}
 \qquad(N\ge0)
\end{equation}
holds there.
\end{corollary}

\begin{proof}
Write $F=\Trunc{F}{P}+E$ and $G=\Trunc{G}{Q}+H$, so that $\ordt E\ge P$ and
$\ordt H\ge Q$, while $\ordt\Trunc{F}{P}\ge a$ and $\ordt\Trunc{G}{Q}\ge b$.
Then
\[
 FG-\bigl(\Trunc{F}{P}\bigr)\bigl(\Trunc{G}{Q}\bigr)
 =E\,\Trunc{G}{Q}+\Trunc{F}{P}\,H+EH,
\]
% ed.: the source bounds the third term by $P+Q$ and asserts $P+Q\ge N$.  That
% ed.: inequality does not follow from $P\ge N-b$ and $Q\ge N-a$: for
% ed.: $a=b=0$ and $N=P=Q=-5$ one has $P+Q=-10<N$.  The repair is to bound
% ed.: $\ordt H$ by $b$ rather than by $Q$, which is legitimate because every
% ed.: block of $H$ is a block of $G$.
whose $\ordt$ is at least $P+b\ge N$, $a+Q\ge N$ and $P+b\ge N$
respectively; for the third term note that $H=G-\Trunc{G}{Q}$ has
$\ordt H\ge b$, every block of $H$ being a block of $G$, so
$\ordt(EH)\ge P+b$.  Hence the two products agree below $t^N$, which is
\eqref{plt:eq:mul-truncated-product}.

For the optimality and the sharpness, assume $N>a+b$, so that $N-1\ge a+b$.
By \cref{plt:thm:mul-finite-determination} the block $[t^{N-1}](FG)$ already
depends on $p_{N-1-b}$ and on $q_{N-1-a}$, so no pair below $(N-b,N-a)$ can
serve for all $F,G$; the following pair of series shows that neither entry
may be lowered by one.  If $N-b-1>a$ put $F=t^a+t^{N-b-1}$; if
$N-b-1=a$ put $F=t^a$.  In both cases take $G=t^b$.  Then
$[t^{N-1}](FG)=1$, whereas $\Trunc{F}{N-b-1}$ omits the block $t^{N-b-1}$,
so $\Trunc{(\Trunc{F}{N-b-1})G}{N}$ has zero block at $t^{N-1}$.  (In the
first case $N-1>a+b$, so the surviving term $t^{a+b}$ does not sit at
level $N-1$.)  The second cutoff is symmetric.

For \eqref{plt:eq:mul-truncated-product-crude}, note that $a,b\ge0$ gives
$N\ge N-b$ and $N\ge N-a$, so $(P,Q)=(N,N)$ is admissible.
\end{proof}

\begin{warningbox}
The crude law \eqref{plt:eq:mul-truncated-product-crude} is \emph{false} on
the Laurent ring.  Take $F=t^{-1}$, $G=1+t$, $N=1$.  Then $FG=t^{-1}+1$, so
$\Trunc{FG}{1}=t^{-1}+1$, whereas
$\Trunc{F}{1}\Trunc{G}{1}=t^{-1}\cdot1=t^{-1}$.  The correct cutoff for $G$
is $N-\ordt F=2$, and indeed $t^{-1}(1+t)=t^{-1}+1$.  Implementations that
``normalize every series to start at $t^0$'' silently commit this error
whenever a genuine principal part is present.
\end{warningbox}

% ed.: S6, Theorem "Closure and finite determination", eq. (mult-error-rule).
% ed.: The source states the conclusion as O(t^{min(p+m_0, q+n_0)}) without
% ed.: any hypothesis relating p,q to the valuations n_0,m_0, and its proof
% ed.: silently bounds ord(E*Gtilde) by p+m_0 and ord(Ftilde*H) by q+n_0,
% ed.: which requires ord(Gtilde) >= m_0 and ord(Ftilde) >= n_0.  Neither
% ed.: follows from the stated hypotheses.  The corrected statement carries
% ed.: the third entry p+q, and the hypotheses p >= n_0, q >= m_0 under
% ed.: which the source's form is recovered.  Note that these hypotheses are
% ed.: sufficient but not necessary: the third entry is redundant exactly
% ed.: when p >= a or q >= b.
\begin{corollary}[Error rule for approximate factors]
\label{plt:cor:mul-error-rule}
Let $F,G\in\PLlau$ with $\ordt F=a$, $\ordt G=b$, and suppose
$\widetilde F,\widetilde G\in\PLlau$ satisfy
\[
 F=\widetilde F+\FormalBigOAt{t^{p}}{t},
 \qquad
 G=\widetilde G+\FormalBigOAt{t^{q}}{t}.
\]
Then
\begin{equation}\label{plt:eq:mul-error-rule}
 FG=\widetilde F\widetilde G
 +\FormalBigOAt{t^{\min(p+b,\;q+a,\;p+q)}}{t}.
\end{equation}
If in addition $p\ge a$ and $q\ge b$ --- that is, if neither approximation is
coarser than the leading order of the quantity it approximates --- then the
third entry is redundant and
\begin{equation}\label{plt:eq:mul-error-rule-normalized}
 FG=\widetilde F\widetilde G
 +\FormalBigOAt{t^{\min(p+b,\;q+a)}}{t}.
\end{equation}
\end{corollary}

\begin{proof}
Write $E=F-\widetilde F$ and $H=G-\widetilde G$, so $\ordt E\ge p$ and
$\ordt H\ge q$.  Expanding,
\begin{equation}\label{plt:eq:mul-error-expansion}
 FG-\widetilde F\widetilde G
 =E\widetilde G+\widetilde FH+EH .
\end{equation}
Now $\widetilde G=G-H$ gives
$\ordt\widetilde G\ge\min(\ordt G,\ordt H)\ge\min(b,q)$, and likewise
$\ordt\widetilde F\ge\min(a,p)$.  Since $\ordt$ is additive on products
(\cref{plt:thm:mul-leading-law} below, or directly
\cref{plt:lem:mul-finite-fiber}) and satisfies
$\ordt(A+B)\ge\min(\ordt A,\ordt B)$,
\[
 \ordt\bigl(FG-\widetilde F\widetilde G\bigr)
 \ \ge\
 \min\bigl(p+\min(b,q),\ \min(a,p)+q,\ p+q\bigr)
 \ =\ \min(p+b,\ q+a,\ p+q),
\]
which is \eqref{plt:eq:mul-error-rule}.  If $p\ge a$ and $q\ge b$, then
$p+q\ge a+q$ and $p+q\ge p+b$, so $p+q\ge\min(p+b,q+a)$ and the third entry
may be dropped.
\end{proof}

\begin{remark}[What \eqref{plt:eq:mul-error-rule} does and does not say]
\label{plt:rmk:mul-not-an-ideal}
The symbol $\FormalBigOAt{t^{N}}{t}$ asserts coefficient agreement below
outer order $N$ and nothing else: no numerical bound, no convergence, no
uniformity in $L$.  It is also not a statement about a quotient ring.  In
$\PLlau$ the element $t$ is a unit, so the ideal generated by $t^N$ is the
whole ring; the subset $t^N\PLpow$ is an additive piece of the
valuation filtration, not an ideal.  Congruences of the form
\eqref{plt:eq:mul-error-rule} must therefore be read as valuation
inequalities $\ordt(F-G)\ge N$, or equivalently as equalities of exclusive
truncations.  Quotient-algebra language becomes legitimate only after the
finite principal part has been removed and one works inside $\PLpow$, where
$(t^N)$ is a genuine ideal.
\end{remark}

The corresponding algorithm is the schoolbook convolution restricted to the
window supplied by \cref{plt:thm:mul-finite-determination}.

\begin{lstlisting}[style=pseudo,caption={Truncated block product},label={plt:alg:mul-truncated}]
# F, G sparse maps  exponent -> coefficient in R;  a = ord(F), b = ord(G)
# returns the blocks of Trunc_{<N}(F*G)
function block_product(F, G, a, b, N):
    C := empty map
    for (n, Pn) in F with a <= n <= N-b-1:
        for (m, Qm) in G with b <= m <= N-n-1:
            C[n+m] := C.get(n+m, 0) + Pn*Qm
    normalize every coefficient of C          # collect in L before zero-testing
    delete the blocks of C that are exactly 0
    return C
\end{lstlisting}

\begin{remark}[Cost, and what actually dominates it]
\label{plt:rmk:mul-cost}
With dense consecutive support the loop performs
$\BigO\bigl((N-a-b)^2\bigr)$ coefficient products; with sparse supports of
sizes $s_F,s_G$ it performs at most $s_Fs_G$.  If the retained blocks have
logarithmic degree $\BigO(d)$, each coefficient product costs
$\BigO(d^2)$ scalar operations by schoolbook polynomial arithmetic, so the
total is $\BigO\bigl((N-a-b)^2d^2\bigr)$; fast convolution in $t$ and fast
polynomial multiplication in $L$ reduce both factors.  These are complexity
$\BigO$'s and carry no analytic meaning.  In practice, at the moderate orders
typical of hand asymptotics, exact coefficient normalization --- expansion,
collection, and gcd reduction if the coefficients are rational functions of
$L$ --- costs more than the convolution itself, and the decisive
implementation choice is to never form a product whose outer order exceeds
the requested cutoff.
\end{remark}

\begin{remark}[A rectangular cutoff is not a truncation]
\label{plt:rmk:mul-rectangular-cutoff}
Discarding all monomials $t^nL^j$ with $|n|+|j|$ or $|j|$ above a threshold
does not respect dominance.  With a cap $|j|\le10$ one deletes $t^2L^{100}$
while retaining $t^3L^{-100}$, although
$t^3L^{-100}\precdom t^2L^{100}$ by an enormous margin.  Bounding
$\degL p_n$ is a second, independent approximation; it is legitimate only
when accompanied by a proved degree bound such as
\cref{plt:prop:mul-affine-profile}.
\end{remark}

\subsection{Leading data and the leading-term law}
\label{plt:sec:mul-leading}

\begin{definition}[Leading data]
\label{plt:def:mul-leading-data}
Let $R$ be a logarithmically graded domain over $\K$ and let
$F=\sum_{n\ge n_0}p_nt^n\in R((t))$ be nonzero with
$\ordt F=\min\SetOf{n:p_n\ne0}$; set $\ordt0=+\infty$.  The
\emph{leading block} of $F$ is the whole coefficient $p_{\ordt F}\in R$, the
\emph{leading monomial} is
\begin{equation}\label{plt:eq:mul-leading-monomial}
 \lm(F)\DefinitionEquals t^{\ordt F}L^{\degL p_{\ordt F}},
\end{equation}
and the \emph{leading scalar} is
$\lc(F)\DefinitionEquals\lc_L\bigl(p_{\ordt F}\bigr)\in\K^\times$.  The
\emph{rank-two valuation} of $F$ is
\begin{equation}\label{plt:eq:mul-rank-two}
 v(F)\DefinitionEquals\bigl(\ordt F,\ -\degL p_{\ordt F}\bigr)
 \in\IntegerNumbers^2,
\end{equation}
with $\IntegerNumbers^2$ lexicographically ordered and $v(0)=+\infty$.  These
are three different objects and none is determined by the coarse integer
$\ordt F$ alone.
\end{definition}

% ed.: the leading data, the rank-two valuation, the log-degree profile, the
% ed.: exclusive truncation law and the summability calculus are all
% ed.: introduced in the ring chapter.  Restating them here without saying so
% ed.: would leave two definitions of the same symbols standing in one volume.
% ed.: This remark records that the restatements are verbatim generalizations
% ed.: over the coefficient ring and agree with the originals.
\begin{remark}[These are the objects of the ring chapter, over a general
coefficient ring]
\label{plt:rmk:mul-leading-restated}
\cref{plt:def:mul-leading-data} introduces no new object.  It repeats
\cref{plt:def:lead-leading-data} and \cref{plt:def:ring-rank-two} with the
three coefficient rings $\K[L]$, $\K(L)$, $\K((L^{-1}))$ replaced by an
arbitrary logarithmically graded domain $R$, so that the leading-term law can
be proved once for all four ambient classes simultaneously; on those four
classes the two readings agree.  The same applies to
\cref{plt:def:mul-profile} against \cref{plt:def:lead-profile} --- whose
codomain is $\NaturalNumbers\cup\SetOf{-\infty}$ because there the blocks are
polynomials, whereas over $\K(L)$ and $\K((L^{-1}))$ the value $\degL p_n$ may
be negative --- and to \cref{plt:cor:mul-truncated-product} against
\cref{plt:prop:trunc-product}, which the corollary strengthens to arbitrary
sufficient cutoffs $(P,Q)$ and supplements with optimality and sharpness.
Nothing below re-proves a ring-chapter result for its own sake; the
restatements exist only to make the general coefficient ring available.
\end{remark}

\begin{lemma}[Dominance is growth]
\label{plt:lem:mul-dominance}
On the positive ray, for integers $n,m$ and $j,k$,
\[
 t^nL^j\succdom t^mL^k
 \quad\Longleftrightarrow\quad
 \lim_{X\to+\infty}\frac{X^{-n}(\log X)^j}{X^{-m}(\log X)^k}=+\infty
 \quad\Longleftrightarrow\quad
 n<m,\ \text{or}\ (n=m\ \text{and}\ j>k).
\]
Consequently, for $F$ in $\PLlau$ or in $\PLit$, $\lm(F)$ is the
$\succdom$-largest element of $\CoefficientSupportOf{F}$.  (For $\PLrat$ the
statement is meaningful only after the expansion of
\cref{plt:rmk:mul-rat-no-support}.)
\end{lemma}

\begin{proof}
% ed.: the first equivalence of the display was left unargued.  It is not a
% ed.: theorem of this section at all: $\succdom$ is defined by the
% ed.: right-hand condition in \cref{plt:def:ring-dominance}, so only the
% ed.: analytic equivalence needs an argument, and that argument restates
% ed.: \cref{plt:lem:ring-dominance-analytic} for integer exponents.
By \cref{plt:def:ring-dominance} the relation $\succdom$ is \emph{defined} by
the right-hand condition, so the first equivalence is a definition and only
the second requires an argument; it is
\cref{plt:lem:ring-dominance-analytic} specialized to integer exponents, and
we repeat the two lines.  The ratio equals $X^{m-n}(\log X)^{j-k}$.  If $m>n$
it tends to $+\infty$
because $\log X=\LittleOAt{X^{\varepsilon}}{X\to+\infty}$ for every
$\varepsilon>0$; if $m=n$ it is $(\log X)^{j-k}$, which tends to $+\infty$
exactly when $j>k$; if $m<n$ it tends to $0$ for the same reason as the first
case.  For the last sentence: a monomial $t^nL^j$ in the support has
$n\ge\ordt F$, with equality attained; among those with $n=\ordt F$ the
largest $j$ with $\kappa_j(p_{\ordt F})\ne0$ is by definition
$\degL p_{\ordt F}$.
\end{proof}

\begin{theorem}[Leading-term law]
\label{plt:thm:mul-leading-law}
Let $R$ be a logarithmically graded domain over $\K$ and let $F,G\in R((t))$
be nonzero.  Then $FG\ne0$ and
\begin{align}
 \ordt(FG)&=\ordt F+\ordt G,
 \label{plt:eq:mul-ord-additive}\\
 \lm(FG)&=\lm(F)\,\lm(G),
 \label{plt:eq:mul-lm-law}\\
 \lc(FG)&=\lc(F)\,\lc(G),
 \label{plt:eq:mul-lc-law}\\
 v(FG)&=v(F)+v(G).
 \label{plt:eq:mul-v-additive}
\end{align}
Moreover $\ordt(F+G)\ge\min(\ordt F,\ordt G)$ and
$v(F+G)\ge\min\bigl(v(F),v(G)\bigr)$ in the lexicographic order.  Each is an
equality unless cancellation occurs at the common level: whole leading blocks
cancel in the first case, top logarithmic coefficients in the second.
\end{theorem}

\begin{proof}
Put $a=\ordt F$, $b=\ordt G$.  By \cref{plt:lem:mul-finite-fiber} the
coefficient $C_N$ of $FG$ vanishes for $N<a+b$, and for $N=a+b$ the index set
$[a,N-b]=\SetOf{a}$ is a single point, so $C_{a+b}=p_aq_b$.  Since $R$ is a
domain (\cref{plt:lem:mul-graded-basics}) and $p_a,q_b\ne0$, we get
$C_{a+b}\ne0$; hence $FG\ne0$ and \eqref{plt:eq:mul-ord-additive} holds.
Applying \eqref{plt:eq:mul-degL-laws} to $C_{a+b}=p_aq_b$ gives
\[
 \degL C_{a+b}=\degL p_a+\degL q_b,
 \qquad
 \lc_L(C_{a+b})=\lc_L(p_a)\lc_L(q_b),
\]
which are exactly \eqref{plt:eq:mul-lm-law} and \eqref{plt:eq:mul-lc-law};
\eqref{plt:eq:mul-v-additive} is the same pair of identities rewritten in the
notation \eqref{plt:eq:mul-rank-two}.

For the sum, let $c=\min(a,b)$.  Every block of $F+G$ of index $<c$ vanishes,
so $\ordt(F+G)\ge c$.  If $a<b$ then $[t^c](F+G)=p_a\ne0$ and the leading
data of $F+G$ and $F$ coincide, so $v(F+G)=v(F)=\min(v(F),v(G))$ --- note
that $a<b$ makes $v(F)<v(G)$ lexicographically.  If $a=b=c$ and
$p_c+q_c\ne0$, then $\ordt(F+G)=c$ and
$\degL(p_c+q_c)\le\max(\degL p_c,\degL q_c)$ by
\eqref{plt:eq:mul-degL-laws}, i.e.
$-\degL(p_c+q_c)\ge\min(-\degL p_c,-\degL q_c)$, which is
$v(F+G)\ge\min(v(F),v(G))$.  If $p_c+q_c=0$, then $\ordt(F+G)>c$ and
$v(F+G)>\min(v(F),v(G))$ outright.
\end{proof}

\begin{corollary}[Integrality and fields]
\label{plt:cor:mul-domain}
Each of $\PLpow$, $\PLlau$, $\PLrat$, $\PLit$ is an integral domain.  The map
$v$ of \eqref{plt:eq:mul-rank-two} is a rank-two valuation on $\PLrat$ and on
$\PLit$, with value group $\IntegerNumbers^2$ lexicographically ordered; on
$\PLlau$ it takes values in the sub-semigroup
$\IntegerNumbers\times\IntegerNumbers_{\le0}$.  The classes $\PLrat$ and
$\PLit$ are fields; $\PLpow$ and $\PLlau$ are not.
\end{corollary}

\begin{proof}
Integrality is \cref{plt:thm:mul-leading-law}, and the valuation axioms are
its last two displays together with $v(F)=+\infty$ iff $F=0$.  Surjectivity
onto $\IntegerNumbers^2$ for $\PLrat$ and $\PLit$ follows from
$v(t^nL^{-j})=(n,j)$, legitimate there because $L^{-j}$ lies in $\K(L)$ and
in $\K((L^{-1}))$; the group $\IntegerNumbers^2$ with the lexicographic order
has exactly one proper nontrivial convex subgroup,
$\SetOf{0}\times\IntegerNumbers$, so the valuation has rank two.  On $\PLlau$
the second coordinate is $-\degL p_{\ordt F}\le0$ because polynomial degrees
are nonnegative.

A Laurent-series ring $R((t))$ over a \emph{field} $R$ is a field, which
covers $\PLrat$ and $\PLit$: given $F\ne0$ with $\nu=\ordt F$ and leading
block $p_\nu\in R^\times$, write $F=p_\nu t^{\nu}(1+U)$ with
$U=\sum_{n>\nu}(p_n/p_\nu)t^{n-\nu}\in tR[[t]]$; the family
$\bigl((-U)^r\bigr)_{r\ge0}$ is summable because $\ordt(-U)^r\ge r$, and by
\cref{plt:lem:mul-summable}(2) its sum $V$ satisfies
\[
 (1+U)V=\sum_{r\ge0}(-U)^r-\sum_{r\ge0}(-U)^{r+1}=(-U)^0=1,
\]
so $F^{-1}=p_\nu^{-1}t^{-\nu}V$.  (\cref{plt:lem:mul-summable} is proved in
\cref{plt:sec:mul-powers} from \cref{plt:thm:mul-leading-law} alone, so there
is no circularity.)  Finally $L\in\PLlau$ is not invertible in
$\PLlau$: if $LH=1$ then \eqref{plt:eq:mul-lm-law} gives
$L\cdot\lm(H)=1$, forcing $\degL$ of the leading block of $H$ to be $-1$,
impossible for a polynomial.
% ed.: the statement asserts that neither $\PLpow$ nor $\PLlau$ is a field,
% ed.: but the argument as given covers only $\PLlau$.  One line closes it.
Since $L\in\PLpow\subseteq\PLlau$, an inverse of $L$ inside $\PLpow$ would in
particular be an inverse inside $\PLlau$; so $L$ is not invertible in
$\PLpow$ either, and neither class is a field.
\end{proof}

\begin{corollary}[Necessity half of the unit criterion]
\label{plt:cor:mul-unit-necessary}
If $F\in\PLlau$ is invertible in $\PLlau$, then its leading block is a
nonzero scalar: $p_{\ordt F}\in\K^\times$.  The same statement holds in
$\PLpow$ with the additional conclusion $\ordt F=0$.
\end{corollary}

\begin{proof}
If $FH=1$ then \eqref{plt:eq:mul-ord-additive} gives
$\ordt F+\ordt H=0$ and \eqref{plt:eq:mul-lm-law} gives
$\degL p_{\ordt F}+\degL h_{\ordt H}=0$.  Both degrees are nonnegative
integers, hence both vanish, so $p_{\ordt F}\in\K^\times$.  In $\PLpow$ both
orders are $\ge0$ and sum to $0$, so both are $0$.
\end{proof}

The converse --- that a scalar leading block suffices --- is a division
statement and is proved with the reciprocal recurrence; it is not a
multiplication fact and is not repeated here.

\subsection{Degree and support bounds}
\label{plt:sec:mul-degree}

\begin{definition}[Log-degree profile]
\label{plt:def:mul-profile}
For $F=\sum_np_nt^n$ in one of the four classes, the \emph{log-degree
profile} is the function $\dprofile{F}\colon\IntegerNumbers\to
\IntegerNumbers\cup\SetOf{-\infty}$,
\begin{equation}\label{plt:eq:mul-profile}
 \dprofile{F}(n)\DefinitionEquals
 \begin{cases}
  \degL p_n,&p_n\ne0,\\
  -\infty,&p_n=0.
 \end{cases}
\end{equation}
The profile is data additional to $\ordt F$; it is not recoverable from it.
For two profiles define the \emph{max-plus convolution}
\begin{equation}\label{plt:eq:mul-maxplus}
 \bigl(\dprofile{F}\ast\dprofile{G}\bigr)(N)
 \DefinitionEquals
 \max_{n+m=N}\bigl(\dprofile{F}(n)+\dprofile{G}(m)\bigr),
\end{equation}
with the conventions $-\infty+c=-\infty$ and $\max\emptyset=-\infty$.  Only
the finitely many pairs of \cref{plt:lem:mul-finite-fiber} can contribute a
value other than $-\infty$, so the maximum is effectively over a finite set.
\end{definition}

\begin{theorem}[Degree bound and exact criterion for attainment]
\label{plt:thm:mul-degree-bound}
Let $F,G\in R((t))$ with $R$ a logarithmically graded domain over $\K$, let
$N\in\IntegerNumbers$, and set
$M\DefinitionEquals(\dprofile{F}\ast\dprofile{G})(N)$.  Let
\[
 P_N\DefinitionEquals
 \SetOf{(n,m):n+m=N,\ p_n\ne0,\ q_m\ne0,\
        \dprofile{F}(n)+\dprofile{G}(m)=M}
\]
be the (finite, possibly empty) set of maximizing pairs.  Then
\begin{equation}\label{plt:eq:mul-degree-bound}
 \dprofile{FG}(N)\le M,
\end{equation}
% ed.: the criterion was stated as an unconditional ``if and only if''.  It
% ed.: fails in the degenerate case $P_N=\emptyset$ (equivalently
% ed.: $M=-\infty$): there $[t^N](FG)=0$, so equality holds although the empty
% ed.: sum \eqref{plt:eq:mul-top-coefficient} is zero.  The criterion is now
% ed.: confined to $M$ finite and the degenerate case is recorded separately.
If $P_N=\emptyset$, equivalently $M=-\infty$, then $[t^N](FG)=0$ and
\eqref{plt:eq:mul-degree-bound} holds with equality.  If $P_N\ne\emptyset$,
so that $M\in\IntegerNumbers$, then the coefficient of $L^{M}$ in the block
$[t^N](FG)$ is exactly
\begin{equation}\label{plt:eq:mul-top-coefficient}
 \kappa_{M}\bigl([t^N](FG)\bigr)
 =\sum_{(n,m)\in P_N}\lc_L(p_n)\,\lc_L(q_m),
\end{equation}
and equality holds in \eqref{plt:eq:mul-degree-bound} if and only if the
sum \eqref{plt:eq:mul-top-coefficient} is nonzero.
\end{theorem}

\begin{proof}
Each summand of $C_N=\sum_{n+m=N}p_nq_m$ has
$\degL(p_nq_m)=\dprofile{F}(n)+\dprofile{G}(m)\le M$ by
\eqref{plt:eq:mul-degL-laws}, with the convention that a vanishing factor
gives $-\infty$.  A finite sum of elements of $\degL\le M$ has
$\degL\le M$, which is \eqref{plt:eq:mul-degree-bound}.  If $M=-\infty$ then
every summand $p_nq_m$ vanishes, so $[t^N](FG)=0$ and both sides of
\eqref{plt:eq:mul-degree-bound} equal $-\infty$; assume from now on
$P_N\ne\emptyset$, so that $M\in\IntegerNumbers$.  By $\K$-linearity of
$\kappa_M$,
\[
 \kappa_M(C_N)=\sum_{n+m=N}\kappa_M(p_nq_m).
\]
If $(n,m)\notin P_N$ then either a factor vanishes or
$\degL(p_nq_m)<M$, and in both cases $\kappa_M(p_nq_m)=0$.  If
$(n,m)\in P_N$ then $\degL(p_nq_m)=M$ and
$\kappa_M(p_nq_m)=\lc_L(p_nq_m)=\lc_L(p_n)\lc_L(q_m)$ by
\cref{plt:lem:mul-graded-basics}.  This proves
\eqref{plt:eq:mul-top-coefficient}, and equality in
\eqref{plt:eq:mul-degree-bound} means precisely $\kappa_M(C_N)\ne0$.
\end{proof}

% ed.: S6, Proposition "Degree and support bounds", asserts equality "whenever
% ed.: the maximizing blocks are nonzero" under the hypothesis that all
% ed.: coefficients are nonnegative over an ordered field.  That hypothesis is
% ed.: far stronger than necessary and the quoted sufficient condition is
% ed.: stated ambiguously (a zero block has degree -infinity and cannot
% ed.: maximize at all).  The criterion below is exact and the positivity
% ed.: hypothesis has been reduced to the leading coefficients of the
% ed.: maximizing blocks only.
\begin{corollary}[Three sufficient conditions for equality]
\label{plt:cor:mul-degree-equality}
In the situation of \cref{plt:thm:mul-degree-bound}, equality holds in
\eqref{plt:eq:mul-degree-bound} in each of the following cases.
\begin{enumerate}
\item $P_N$ is a singleton.
\item $N=\ordt F+\ordt G$ (both factors nonzero); then $P_N$ is the
singleton $\SetOf{(\ordt F,\ordt G)}$ and one recovers
\eqref{plt:eq:mul-lm-law}.
% ed.: condition (3) as first drafted omitted $P_N\ne\emptyset$; over an
% ed.: ordered field its hypotheses are satisfied vacuously when $P_N$ is
% ed.: empty, and the conclusion would then be the false assertion that the
% ed.: sum \eqref{plt:eq:mul-top-coefficient} is nonzero.
\item $P_N\ne\emptyset$, $\K$ is an ordered field, and $\lc_L(p_n)>0$,
$\lc_L(q_m)>0$ for every $(n,m)\in P_N$.
\end{enumerate}
\end{corollary}

\begin{proof}
(1) The sum \eqref{plt:eq:mul-top-coefficient} is a single product of two
elements of $\K^\times$, hence nonzero.  (2) By
\cref{plt:lem:mul-finite-fiber} the only admissible pair at
$N=\ordt F+\ordt G$ is $(\ordt F,\ordt G)$, and both blocks are nonzero, so
$P_N$ is that singleton; apply (1).  (3) The sum is a nonempty sum of
strictly positive elements of an ordered field.
\end{proof}

\begin{example}[The bound is attained, at every order and by every pair]
\label{plt:ex:mul-degree-sharp}
Over $\K=\RationalNumbers$ let
\[
 F=G=\sum_{n\ge0}L^nt^n\in\PLpow,
 \qquad\text{so}\qquad
 \dprofile{F}(n)=\dprofile{G}(n)=n .
\]
Then $(\dprofile{F}\ast\dprofile{G})(N)=\max_{n+m=N}(n+m)=N$, and
\emph{every} one of the $N+1$ admissible pairs is maximizing, so $P_N$ is as
large as it can be and \cref{plt:cor:mul-degree-equality}(1) does not apply.
Nevertheless
\[
 [t^N](FG)=\sum_{n=0}^{N}L^n\cdot L^{N-n}=(N+1)L^N,
\]
so $\dprofile{FG}(N)=N$ for every $N$: the bound
\eqref{plt:eq:mul-degree-bound} is attained at every order.  Formula
\eqref{plt:eq:mul-top-coefficient} predicts exactly the top coefficient
$N+1=\sum_{(n,m)\in P_N}1\cdot1$.
\end{example}

\begin{warningbox}
\cref{plt:ex:mul-degree-sharp} is the place where characteristic zero is
used.  Over a field of characteristic $p>0$ the same computation gives
$[t^{p-1}](FG)=p\,L^{p-1}=0$, so the block vanishes entirely and
$\dprofile{FG}(p-1)=-\infty$ although the bound is $p-1$.  Every statement in
this section is asserted over a field of characteristic zero, and this is not
a formality.
\end{warningbox}

\begin{example}[The bound can fail by an arbitrary amount]
\label{plt:ex:mul-degree-drop}
Let $F=1+L^{d}t$ and $G=1-L^{d}t$ in $\PLpow$, $d\ge1$.  Then
$(\dprofile{F}\ast\dprofile{G})(1)=\max(0+d,\ d+0)=d$, while
$[t^1](FG)=-L^{d}+L^{d}=0$, so $\dprofile{FG}(1)=-\infty$.  Here
$P_1=\SetOf{(0,1),(1,0)}$ and the sum
\eqref{plt:eq:mul-top-coefficient} is $1\cdot(-1)+1\cdot1=0$, as it must be.
A drop by exactly one occurs for $F=1+(L^2+L)t$, $G=1-(L^2-L)t$, where the
bound at $N=1$ is $2$ while $[t^1](FG)=2L$ has degree $1$: here the top
coefficients cancel but the next ones do not.
\end{example}

\begin{proposition}[Affine profiles and the profile subalgebra]
\label{plt:prop:mul-affine-profile}
Let $\alpha,\alpha'\ge0$ and $\beta,\gamma\in\RealNumbers$.  Let
$F,G\in\PLlau$ be nonzero with $a=\ordt F$, $b=\ordt G$ and
\[
 \dprofile{F}(n)\le\alpha n+\beta\ \ (n\ge a),
 \qquad
 \dprofile{G}(m)\le\alpha'm+\gamma\ \ (m\ge b).
\]
Put $\mu=\max(\alpha,\alpha')$.  Then for all $N$,
\begin{equation}\label{plt:eq:mul-affine-bound}
 \dprofile{FG}(N)\le\mu N+\beta+\gamma-(\mu-\alpha)a-(\mu-\alpha')b .
\end{equation}
In particular, if $a,b\ge0$ then $\dprofile{FG}(N)\le\mu N+\beta+\gamma$.
Consequently, for each $\alpha\ge0$ the set
\[
 \mathfrak A_\alpha
 \DefinitionEquals
 \SetOf{F\in\PLlau:\ \exists\beta\in\RealNumbers,\
        \dprofile{F}(n)\le\alpha n+\beta\ \text{for all }n}
\]
is a $\K$-subspace with
$\mathfrak A_\alpha\cdot\mathfrak A_{\alpha'}\subseteq
\mathfrak A_{\max(\alpha,\alpha')}$, and
$\mathfrak A\DefinitionEquals\bigcup_{\alpha\ge0}\mathfrak A_\alpha$ is a
$\K$-subalgebra of $\PLlau$ containing $\K[L][t,t^{-1}]$.
\end{proposition}

\begin{proof}
Fix $N$ and $n+m=N$ with $n\ge a$, $m\ge b$.  Then
\[
 \alpha n+\alpha'm
 =\mu(n+m)-(\mu-\alpha)n-(\mu-\alpha')m
 \le\mu N-(\mu-\alpha)a-(\mu-\alpha')b,
\]
because $\mu-\alpha\ge0$, $\mu-\alpha'\ge0$, $n\ge a$ and $m\ge b$.  Adding
$\beta+\gamma$ and taking the maximum over the finite index set gives
\eqref{plt:eq:mul-affine-bound} via \eqref{plt:eq:mul-degree-bound}.  If
$a,b\ge0$ then the two subtracted terms are $\ge0$, so they may be dropped.

For the algebra statements: $\mathfrak A_\alpha$ is a subspace because
$\dprofile{F+G}(n)\le\max(\dprofile{F}(n),\dprofile{G}(n))
\le\alpha n+\max(\beta,\gamma)$ and scalars do not change profiles.  The
product statement is \eqref{plt:eq:mul-affine-bound}, whose right-hand side is
$\mu N+\text{const}$.
% ed.: the inclusion $\mathfrak A_\alpha\subseteq\mathfrak A_{\alpha''}$ for
% ed.: $\alpha\le\alpha''$ was used without proof.  On the Laurent ring it is
% ed.: not immediate: for $n<0$ a larger slope gives a \emph{smaller} bound,
% ed.: so the constant has to be increased, and that is possible only because
% ed.: the support is bounded below.  The argument is supplied here.
For $\alpha\le\alpha''$ one has $\mathfrak A_\alpha\subseteq\mathfrak
A_{\alpha''}$.  Indeed, let $F\ne0$ satisfy $\dprofile{F}(n)\le\alpha n+\beta$
for all $n$, and put $a=\ordt F$.  For $n\ge0$ we have
$\alpha n+\beta\le\alpha''n+\beta$; for $a\le n<0$,
\[
 \alpha n+\beta=\alpha''n+(\alpha''-\alpha)(-n)+\beta
 \le\alpha''n+\beta+(\alpha''-\alpha)\max(0,-a);
\]
and $\dprofile{F}(n)=-\infty$ for $n<a$.  Hence
$\beta''=\beta+(\alpha''-\alpha)\max(0,-a)$ witnesses
$F\in\mathfrak A_{\alpha''}$, and $0\in\mathfrak A_{\alpha''}$ trivially.
Consequently the union $\mathfrak A$ is closed under
addition and multiplication and contains $1$; every Laurent polynomial
$\sum_{n=n_1}^{n_2}p_n(L)t^n$ lies in $\mathfrak A_0$ with
$\beta=\max_n\degL p_n$.
\end{proof}

\begin{remark}[$\mathfrak A$ is not $t$-adically closed]
\label{plt:rmk:mul-profile-not-closed}
The affine-profile condition is preserved by finite sums and products but not
by $t$-adic limits.  Take $F_r=t^rL^{r^2}$ for $r\ge1$; each $F_r$ lies in
$\mathfrak A_0$, and $\ordt F_r=r\to\infty$, so $(F_r)$ is summable, yet the
sum
$S=\sum_{r\ge1}t^rL^{r^2}$ has $\dprofile{S}(r)=r^2$, which is bounded by no
affine function.  Hence $S\notin\mathfrak A$.  A degree bound must therefore
be re-derived, not inherited, whenever an infinite construction is performed.
\end{remark}

\begin{proposition}[Profile of a power]
\label{plt:prop:mul-power-profile}
Let $F\in\PLpow$ with $\dprofile{F}(n)\le\alpha n+\beta$ for all $n\ge0$,
where $\alpha\ge0$ and $\beta\ge0$.  Then for every integer $k\ge0$,
\begin{equation}\label{plt:eq:mul-power-profile}
 \dprofile{F^k}(N)\le\alpha N+k\beta
 \qquad(N\ge0).
\end{equation}
\end{proposition}

\begin{proof}
Induction on $k$.  For $k=0$ we have $F^0=1$, so $\dprofile{F^0}(0)=0\le0$
and $\dprofile{F^0}(N)=-\infty$ for $N>0$.  Assume the bound for $k$.  By
\eqref{plt:eq:mul-degree-bound} applied to $F^{k+1}=F^k\cdot F$,
\[
 \dprofile{F^{k+1}}(N)
 \le\max_{n+m=N,\ n,m\ge0}\bigl((\alpha n+k\beta)+(\alpha m+\beta)\bigr)
 =\alpha N+(k+1)\beta .
\]
\end{proof}

\begin{remark}[The Lambert profile, corrected]
\label{plt:rmk:mul-lambert-profile}
Two of the sources record the degree law of the canonical Lambert family as
$\deg\LamP{n}=n$ without restriction, and use it to assert that products of
Lambert-type expansions have profile $\le N$.  The law
$\deg\LamP{n}=n$ holds only for $n\ge1$: the normalization
$\LamP{0}(u)=-u$ has degree $1$, not $0$.  The correct hypothesis for a
series $F=\sum_{n\ge0}\LamP{n}(L)t^n$ is therefore
$\dprofile{F}(n)\le n+1$ for all $n\ge0$, i.e. $\alpha=1$, $\beta=1$ in
\cref{plt:prop:mul-affine-profile}; a product of two such series has
$\dprofile{FG}(N)\le N+2$, and a $k$-th power has
$\dprofile{F^k}(N)\le N+k$ by \cref{plt:prop:mul-power-profile}.  The
unrestricted claim $\dprofile{FG}(N)\le N$ is false already at $N=0$, where
$[t^0](FG)=\LamP{0}(L)^2=L^2$ has degree $2$.
\end{remark}

\subsection{Powers and binomial powers}
\label{plt:sec:mul-powers}

Infinite sums enter through powers, so we first record the exact sense in
which they are legitimate.

\begin{definition}[Summable family]
\label{plt:def:mul-summable}
Let $R$ be a logarithmically graded domain over $\K$.  A family
$(S_i)_{i\in I}$ in $R((t))$ is \emph{summable} if for every
$N\in\IntegerNumbers$ the set $\SetOf{i\in I:\ordt S_i<N}$ is finite.  Its
sum is defined blockwise:
$[t^n]\sum_{i}S_i\DefinitionEquals\sum_i[t^n]S_i$, a finite sum in the
coefficient ring.
\end{definition}

\begin{lemma}[Summable families are well behaved]
\label{plt:lem:mul-summable}
Let $(S_i)_{i\in I}$ and $(T_j)_{j\in J}$ be summable families in $R((t))$.
\begin{enumerate}
\item $\inf_i\ordt S_i>-\infty$, so $\sum_iS_i$ is a genuine element of
$R((t))$ (its support is bounded below).
\item The family $(S_iT_j)_{(i,j)\in I\times J}$ is summable and
$\bigl(\sum_iS_i\bigr)\bigl(\sum_jT_j\bigr)=\sum_{(i,j)}S_iT_j$.
\end{enumerate}
\end{lemma}

\begin{proof}
(1) Taking $N=0$ in \cref{plt:def:mul-summable}, only finitely many $i$ have
$\ordt S_i<0$, and each of those has $\ordt S_i\in\IntegerNumbers$ (a member
of that finite set is in particular nonzero).  Let $\lambda$ be the minimum of
that finite set of integers, or $\lambda=0$ if the set is empty.  Then
$\ordt S_i\ge\min(0,\lambda)$ for all $i$, and every block of
$\sum_iS_i$ below that bound vanishes.

% ed.: the source calls the two infima ``finite by (1)''; part~(1) gives only
% ed.: $>-\infty$, and the value $+\infty$ does occur, namely when every
% ed.: member of the family is zero.  That case is disposed of separately.
(2) Let $\alpha=\inf_i\ordt S_i$ and $\beta=\inf_j\ordt T_j$, both
$>-\infty$ by (1).  If $\alpha=+\infty$ or $\beta=+\infty$ then every $S_i$,
respectively every $T_j$, is zero, so every $S_iT_j$ is zero and both
assertions are trivial; assume therefore $\alpha,\beta\in\IntegerNumbers$.
Fix $N$.  If $\ordt(S_iT_j)=\ordt S_i+\ordt T_j<N$ then
$\ordt S_i<N-\beta$ and $\ordt T_j<N-\alpha$, and each of those conditions
holds for only finitely many indices; so the family is summable.  For the
identity, fix $n$ and compute in the coefficient ring, all sums having only
finitely many nonzero terms:
\[
 [t^n]\Bigl(\sum_iS_i\sum_jT_j\Bigr)
 =\sum_{p+q=n}\Bigl(\sum_i[t^p]S_i\Bigr)\Bigl(\sum_j[t^q]T_j\Bigr)
 =\sum_{i,j}\sum_{p+q=n}[t^p]S_i\,[t^q]T_j
 =\sum_{i,j}[t^n](S_iT_j).
\]
\end{proof}

\begin{remark}
\label{plt:rmk:mul-summable-vs-cauchy}
Part (1) is worth isolating.  For a \emph{sequence} in a Laurent ring,
$t$-adic Cauchyness alone does not guarantee a limit with a finite principal
part; a common eventual lower support bound must be imposed.  For a
\emph{summable family} no such extra hypothesis is needed, because the
finiteness condition at $N=0$ already forces a common lower bound.
\end{remark}

\begin{proposition}[Coefficients of an integer power]
\label{plt:prop:mul-integer-power}
Let $F=\sum_{n\ge0}p_nt^n\in\PLpow$ and $k\ge0$.  Then
\begin{equation}\label{plt:eq:mul-power-coefficient}
 [t^N]F^k
 =\sum_{\substack{n_1,\ldots,n_k\ge0\\ n_1+\cdots+n_k=N}}
   p_{n_1}\cdots p_{n_k}
 \DefinitionEquals C_{N,k}(p_0,p_1,\ldots),
\end{equation}
with $C_{0,0}=1$ and $C_{N,0}=0$ for $N>0$.  These \emph{convolution
polynomials} satisfy
\begin{equation}\label{plt:eq:mul-Cmr-recurrence}
 C_{N,k+1}=\sum_{n=0}^{N}p_n\,C_{N-n,k},
\end{equation}
which computes the whole triangular table in $\BigO(N^2k)$ coefficient
products.  If $p_0=0$, then $C_{N,k}=0$ for $N<k$, and, writing
$F=\sum_{n\ge1}p_nt^n$,
\begin{equation}\label{plt:eq:mul-bell}
 [t^N]F^k
 =\frac{k!}{N!}\,
 \ExponentialPartialBellPolynomial{N}{k}
   \bigl(1!\,p_1,\,2!\,p_2,\ldots,(N-k+1)!\,p_{N-k+1}\bigr).
\end{equation}
\end{proposition}

\begin{proof}
Formula \eqref{plt:eq:mul-power-coefficient} follows by induction on $k$ from
\cref{plt:def:mul-cauchy}: the case $k=0$ is the convention, and
\[
 [t^N]F^{k+1}=\sum_{n+m=N}p_n\,[t^m]F^k
 =\sum_{n=0}^{N}p_n\!\!\!\!\sum_{\substack{n_1,\ldots,n_k\ge0\\
   n_1+\cdots+n_k=N-n}}\!\!\!\!
   p_{n_1}\cdots p_{n_k},
\]
which is both \eqref{plt:eq:mul-Cmr-recurrence} and the composition sum for
$k+1$.  All sums are finite because every $n_i\ge0$ and they sum to $N$.  If
$p_0=0$, every $n_i$ in a nonvanishing term is $\ge1$, so $N\ge k$.

For \eqref{plt:eq:mul-bell}, recall the defining generating identity of the
exponential partial Bell polynomials \cite{comtet-book,comtet1970}: in
$R[[z]]$ over a commutative $\RationalNumbers$-algebra $R$,
\[
 \frac1{k!}\Bigl(\sum_{j\ge1}x_j\frac{z^j}{j!}\Bigr)^{k}
 =\sum_{N\ge k}\ExponentialPartialBellPolynomial{N}{k}(x_1,x_2,\ldots)
   \frac{z^N}{N!}.
\]
Apply it with $R=\K[L]$ (a $\RationalNumbers$-algebra, as
$\operatorname{char}\K=0$), $z=t$ and $x_j=j!\,p_j$, so that
$\sum_{j\ge1}x_jz^j/j!=F$.  Comparing coefficients of $t^N$ gives
$\frac1{k!}[t^N]F^k=\frac1{N!}
\ExponentialPartialBellPolynomial{N}{k}(1!p_1,2!p_2,\ldots)$, which is
\eqref{plt:eq:mul-bell}; only the arguments $x_1,\ldots,x_{N-k+1}$ occur
because a term of $\ExponentialPartialBellPolynomial{N}{k}$ is a product of
$k$ variables with index sum $N$, so no index can exceed $N-k+1$.
\end{proof}

\begin{theorem}[Binomial powers and their finite determination]
\label{plt:thm:mul-binomial}
Let $U\in\PLpow$ with $\ordt U\ge e$ for an integer $e\ge1$, write
$u_n=[t^n]U$, and let $\sigma\in\K$.  Put
$\binom{\sigma}{r}=\sigma(\sigma-1)\cdots(\sigma-r+1)/r!\in\K$, with
$\binom{\sigma}{0}=1$.  Then:
\begin{enumerate}
\item The family $\bigl(\binom{\sigma}{r}U^r\bigr)_{r\ge0}$ is summable, so
\begin{equation}\label{plt:eq:mul-binomial-def}
 (1+U)^\sigma\DefinitionEquals\sum_{r\ge0}\binom{\sigma}{r}U^r
\end{equation}
is a well-defined element of $\PLpow$ with constant block $1$.
\item \emph{(Finite determination.)}  For $N\ge0$,
\begin{equation}\label{plt:eq:mul-binomial-finite}
 [t^N](1+U)^\sigma
 =\sum_{r=0}^{\Floor{N/e}}\binom{\sigma}{r}\,[t^N]U^r,
\end{equation}
and this depends only on the blocks $u_e,u_{e+1},\ldots,u_N$, that is, only
on $\Trunc{U}{N+1}$.  The range of $r$ is exact: for $U=t^e$ with $e\mid N$
the term $r=N/e$ contributes $\binom{\sigma}{N/e}$, nonzero whenever
$\sigma\notin\SetOf{0,1,\ldots,N/e-1}$.  The top block enters exactly
linearly: for $N\ge1$,
\begin{equation}\label{plt:eq:mul-binomial-top}
 [t^N](1+U)^\sigma=\sigma u_N+\Psi_N(u_e,\ldots,u_{N-e}),
\end{equation}
where $\Psi_N$ is a universal polynomial with coefficients in
$\RationalNumbers[\sigma]$, with the convention that $\Psi_N=0$ when the
argument list is empty.  Hence for $\sigma\ne0$ the block $u_N$ cannot be
omitted from the dependency range.
\item \emph{(Exponent law.)}  For $\sigma,\tau\in\K$,
\begin{equation}\label{plt:eq:mul-binomial-group}
 (1+U)^\sigma(1+U)^\tau=(1+U)^{\sigma+\tau},
 \qquad (1+U)^0=1,
\end{equation}
and for $k\in\NaturalNumbers$ the series $(1+U)^k$ agrees with the $k$-fold
product.  Consequently $(1+U)^\sigma$ is a unit of $\PLpow$ with inverse
$(1+U)^{-\sigma}$, and $\sigma\mapsto(1+U)^\sigma$ is a homomorphism from
$(\K,+)$ to the units of $\PLpow$.
\item \emph{(Profile.)}  If $\dprofile{U}(n)\le\alpha n+\beta$ for all
$n\ge e$, with $\alpha\ge0$, then, writing $\beta^+=\max(\beta,0)$,
\begin{equation}\label{plt:eq:mul-binomial-profile}
 \dprofile{(1+U)^\sigma}(N)\le\Bigl(\alpha+\frac{\beta^+}{e}\Bigr)N
 \qquad(N\ge0).
\end{equation}
In particular $\mathfrak A\cap\PLpow$ is stable under
$U\mapsto(1+U)^\sigma$ whenever $\ordt U\ge1$.
\end{enumerate}
\end{theorem}

\begin{proof}
(1) By \eqref{plt:eq:mul-ord-additive}, $\ordt(U^r)\ge re$, which tends to
$+\infty$; so for each $N$ only the finitely many $r\le N/e$ can have
$\ordt(\binom{\sigma}{r}U^r)<N$.  \cref{plt:lem:mul-summable}(1) gives a
well-defined sum, and $\ordt U\ge e\ge1$ makes the $r=0$ term the only
contributor to $t^0$.

(2) Blockwise, $[t^N](1+U)^\sigma=\sum_{r\ge0}\binom{\sigma}{r}[t^N]U^r$, and
$[t^N]U^r=0$ once $re>N$; this is \eqref{plt:eq:mul-binomial-finite}.  By
\cref{plt:thm:mul-finite-determination} applied $r-1$ times, $[t^N]U^r$
depends only on the blocks $u_n$ with $e\le n\le N-(r-1)e$, and for $r\ge1$
this range is contained in $[e,N]$.  For the sharpness of the $r$-range take
$U=t^e$: then $U^r=t^{re}$ and the only contribution to $t^N$ comes from
$r=N/e$, with value $\binom{\sigma}{N/e}$, which vanishes exactly when
$\sigma$ is one of $0,1,\ldots,N/e-1$.  For
\eqref{plt:eq:mul-binomial-top}, the term $r=0$ contributes nothing to
$t^N$ for $N\ge1$, the term $r=1$ contributes $\sigma u_N$, and each term
with $r\ge2$ involves only $u_n$ with $n\le N-(r-1)e\le N-e$.

(3) By \cref{plt:lem:mul-summable}(2),
\[
 (1+U)^\sigma(1+U)^\tau
 =\sum_{r,s\ge0}\binom{\sigma}{r}\binom{\tau}{s}U^{r+s}
 =\sum_{m\ge0}\Bigl(\sum_{r+s=m}\binom{\sigma}{r}\binom{\tau}{s}\Bigr)U^m,
\]
the regrouping being legitimate because for each fixed $t$-block only
finitely many pairs $(r,s)$ contribute.  Vandermonde's convolution
\cite{graham-knuth-patashnik}
\[
 \sum_{r+s=m}\binom{\sigma}{r}\binom{\tau}{s}=\binom{\sigma+\tau}{m}
\]
holds for all $\sigma,\tau\in\K$: both sides are polynomial expressions in
$(\sigma,\tau)$ with coefficients in $\RationalNumbers$; they agree for all
pairs of nonnegative integers by comparing coefficients in
$(1+z)^{\sigma}(1+z)^{\tau}=(1+z)^{\sigma+\tau}$ in
$\RationalNumbers[z]$; a polynomial in two variables over the infinite domain
$\RationalNumbers$ vanishing on $\NaturalNumbers^2$ is the zero polynomial;
and $\RationalNumbers\subseteq\K$ because $\operatorname{char}\K=0$, so the
identity persists in $\K$.  This proves
\eqref{plt:eq:mul-binomial-group}; $(1+U)^0=1$ is immediate.  For
$k\in\NaturalNumbers$ we have $\binom{k}{r}=0$ for $r>k$, so
\eqref{plt:eq:mul-binomial-def} is the finite binomial expansion of the
$k$-fold product in the commutative ring $\PLpow$.  Invertibility follows by
taking $\tau=-\sigma$.

(4) By \eqref{plt:eq:mul-binomial-finite} it suffices to bound
$\degL[t^N]U^r$ for $1\le r\le\Floor{N/e}$ (the term $r=0$ contributes only
at $N=0$, where the bound reads $0\le0$).  By
\eqref{plt:eq:mul-power-coefficient} and \eqref{plt:eq:mul-degL-laws},
\[
 \degL[t^N]U^r
 \le\max_{\substack{n_1+\cdots+n_r=N\\ n_i\ge e}}
    \sum_{i=1}^{r}(\alpha n_i+\beta)
 =\alpha N+r\beta
 \le\alpha N+\Floor{N/e}\beta^+
 \le\Bigl(\alpha+\frac{\beta^+}{e}\Bigr)N,
\]
where the second inequality uses $r\le\Floor{N/e}$ when $\beta\ge0$ and
$r\beta\le0\le\Floor{N/e}\beta^+$ when $\beta<0$.  Taking the maximum over
$r$ gives \eqref{plt:eq:mul-binomial-profile}.
\end{proof}

\begin{corollary}[Formal roots]
\label{plt:cor:mul-roots}
With $U$ as above and $p,q\in\IntegerNumbers$, $q\ge1$,
\[
 \bigl((1+U)^{p/q}\bigr)^{q}=(1+U)^{p},
\]
the right-hand side being the ordinary $p$-th power (or the reciprocal of the
$|p|$-th power if $p<0$).  In particular $(1+U)^{1/q}$ is a $q$-th root of
$1+U$ in $\PLpow$, and it is the unique one with constant block $1$.
\end{corollary}

\begin{proof}
Iterating \eqref{plt:eq:mul-binomial-group} $q-1$ times gives
$((1+U)^{p/q})^q=(1+U)^{qp/q}=(1+U)^p$, and
\cref{plt:thm:mul-binomial}(3) identifies the right-hand side with the
ordinary power.  For uniqueness, suppose $V,W\in\PLpow$ have constant block
$1$ and $V^q=W^q=1+U$.  Then $(V-W)\sum_{i=0}^{q-1}V^iW^{q-1-i}=0$; the
second factor has constant block $q\ne0$ (characteristic zero), hence is
nonzero, and $\PLpow$ is a domain by \cref{plt:cor:mul-domain}, so $V=W$.
\end{proof}

\begin{remark}[Infinite products]
\label{plt:rmk:mul-infinite-products}
If $(U_r)_{r\ge1}$ lies in $\PLpow$ with $\ordt U_r\to+\infty$, the partial
products $P_R=\prod_{r\le R}(1+U_r)$ converge $t$-adically.  Indeed, given
$N$ choose $R_N$ with $\ordt U_r\ge N$ for $r>R_N$; for $R>R_N$,
\[
 P_R-P_{R_N}=P_{R_N}\Bigl(\prod_{R_N<r\le R}(1+U_r)-1\Bigr),
\]
and expanding the finite product shows that the bracket is a sum of products
of at least one $U_r$ with $r>R_N$, hence of $\ordt\ge N$; since
$\ordt P_{R_N}\ge0$ we get $\ordt(P_R-P_{R_N})\ge N$.  Thus $(P_R)$ is
$t$-adically Cauchy in the complete ring $\PLpow$ and
$\prod_{r\ge1}(1+U_r)$ is well defined.  On $\PLlau$ the same argument
applies once one notes that $\ordt U_r\to+\infty$ forces all but finitely
many factors to lie in $\PLpow$, so the partial products have a common lower
support bound.
\end{remark}

\subsection{Three worked computations}
\label{plt:sec:mul-worked}

\begin{example}[A Laurent product through four blocks, with degree
bookkeeping]
\label{plt:ex:mul-worked-laurent}
Over $\K=\RationalNumbers$ take
\begin{align*}
 F&=3t^{-2}+(L-1)t^{-1}+(L^2+2)+2L\,t+(L^2-L)t^2
   +\FormalBigOAt{t^{3}}{t},\\
 G&=t^{-1}+(2-L)+(L+1)t+L^2t^2+\FormalBigOAt{t^{3}}{t}.
\end{align*}
Here $\ordt F=-2$ and $\ordt G=-1$; $F$ is known through $\Trunc{F}{3}$ and
$G$ through $\Trunc{G}{3}$.  By the precision rule the product is known
through
\[
 \Trunc{FG}{\min(3+(-1),\,3+(-2))}=\Trunc{FG}{1},
\]
that is, through the block $t^{0}$ inclusive --- four blocks in all, starting
at $t^{-3}$.  The extra block $[t^2]F$ is wasted: it can influence only
$[t^{1}](FG)$ and beyond, for which $[t^{3}]G$ would also be needed.

Writing $A_n=[t^n]F$ and $B_m=[t^m]G$, the convolution
\eqref{plt:eq:mul-cauchy} gives
\begin{align*}
 C_{-3}&=A_{-2}B_{-1}=3,\\
 C_{-2}&=A_{-2}B_{0}+A_{-1}B_{-1}
        =3(2-L)+(L-1)=5-2L,\\
 C_{-1}&=A_{-2}B_{1}+A_{-1}B_{0}+A_{0}B_{-1}\\
       &=3(L+1)+(-L^2+3L-2)+(L^2+2)=6L+3,\\
 C_{0}&=A_{-2}B_{2}+A_{-1}B_{1}+A_{0}B_{0}+A_{1}B_{-1}\\
      &=3L^2+(L^2-1)+(-L^3+2L^2-2L+4)+2L
      =-L^3+6L^2+3 .
\end{align*}
Hence
\begin{equation}\label{plt:eq:mul-worked-laurent}
 FG=3t^{-3}+(5-2L)t^{-2}+(6L+3)t^{-1}+(-L^3+6L^2+3)
   +\FormalBigOAt{t^{1}}{t}.
\end{equation}

Leading data: $\lm(F)=t^{-2}$ with $\lc(F)=3$, $\lm(G)=t^{-1}$ with
$\lc(G)=1$; \cref{plt:thm:mul-leading-law} predicts $\lm(FG)=t^{-3}$ and
$\lc(FG)=3$, confirmed by $C_{-3}=3$.

The degree bookkeeping is instructive.  The profiles are
$\dprofile{F}(-2)=0$, $\dprofile{F}(-1)=1$, $\dprofile{F}(0)=2$,
$\dprofile{F}(1)=1$ and $\dprofile{G}(-1)=0$, $\dprofile{G}(0)=1$,
$\dprofile{G}(1)=1$, $\dprofile{G}(2)=2$.  Then:

\begin{center}
\begin{tabular}{@{}ccccc@{}}
\hline
$N$ & maximizing pairs $P_N$ & bound $M$ & top coefficient
 \eqref{plt:eq:mul-top-coefficient} & $\dprofile{FG}(N)$ \\
\hline
$-3$ & $(-2,-1)$ & $0$ & $3\cdot1=3$ & $0$ \\
$-2$ & $(-2,0),\ (-1,-1)$ & $1$ & $3(-1)+1\cdot1=-2$ & $1$ \\
$-1$ & $(-1,0),\ (0,-1)$ & $2$ & $1\cdot(-1)+1\cdot1=0$ & $1<2$ \\
$0$  & $(0,0)$ & $3$ & $1\cdot(-1)=-1$ & $3$ \\
\hline
\end{tabular}
\end{center}

At $N=-2$ two pairs compete and their contributions $-3$ and $+1$ do not
cancel, so the bound is attained, with top coefficient $-2$: indeed
$C_{-2}=-2L+5$.  At $N=-1$ two pairs compete and cancel exactly, so the
degree drops from the bound $2$ to the actual value $1$: indeed
$C_{-1}=6L+3$.  At $N=0$ the maximizing pair is unique, so
\cref{plt:cor:mul-degree-equality}(1) forces attainment, with top
coefficient $-1$: indeed $C_0=-L^3+\cdots$.  This is exactly the phenomenon
\cref{plt:thm:mul-degree-bound} isolates; no inspection of the finished
polynomial is needed to predict it.
\end{example}

\begin{example}[A square root through four blocks, verified by squaring]
\label{plt:ex:mul-worked-sqrt}
Let
\[
 U=2Lt+(L^2+2)t^2+4L\,t^3+\FormalBigOAt{t^{4}}{t}\in t\PLpow,
\]
so $e=1$ in \cref{plt:thm:mul-binomial}, and take $\sigma=\tfrac12$, with
\[
 \binom{1/2}{0}=1,\quad
 \binom{1/2}{1}=\tfrac12,\quad
 \binom{1/2}{2}=-\tfrac18,\quad
 \binom{1/2}{3}=\tfrac1{16}.
\]
By \eqref{plt:eq:mul-binomial-finite} only $r\le N$ contributes to $t^N$, so
through $t^3$ we need
\[
 U^2=4L^2t^2+(4L^3+8L)t^3+\FormalBigOAt{t^{4}}{t},
 \qquad
 U^3=8L^3t^3+\FormalBigOAt{t^{4}}{t}.
\]
Block by block:
\begin{align*}
 [t^1](1+U)^{1/2}&=\tfrac12(2L)=L,\\
 [t^2](1+U)^{1/2}&=\tfrac12(L^2+2)-\tfrac18(4L^2)
   =\tfrac{L^2}2+1-\tfrac{L^2}2=1,\\
 [t^3](1+U)^{1/2}&=\tfrac12(4L)-\tfrac18(4L^3+8L)+\tfrac1{16}(8L^3)
   =2L-\tfrac{L^3}2-L+\tfrac{L^3}2=L .
\end{align*}
Hence
\begin{equation}\label{plt:eq:mul-worked-sqrt}
 (1+U)^{1/2}=1+L\,t+t^2+L\,t^3+\FormalBigOAt{t^{4}}{t}.
\end{equation}
Two logarithmic cancellations occur, one at each of $t^2$ and $t^3$; the
degree bound \eqref{plt:eq:mul-binomial-profile} with $\alpha=1$,
$\beta^{+}=0$, $e=1$ predicts only $\dprofile{(1+U)^{1/2}}(N)\le N$, and the
actual
profile $0,1,0,1$ is strictly smaller at $N=2,3$.  A degree bound is an upper
bound; it is not a prediction.

\emph{Verification.}  \cref{plt:cor:mul-roots} says the square of
\eqref{plt:eq:mul-worked-sqrt} must reproduce $1+U$ through $t^3$.  Squaring
by \eqref{plt:eq:mul-cauchy}:
\[
 [t^0]=1,\quad
 [t^1]=2L,\quad
 [t^2]=2\cdot1+L^2=L^2+2,\quad
 [t^3]=2L+2L=4L,
\]
which are exactly the blocks $1,u_1,u_2,u_3$ of $1+U$.  Note the precision
accounting: squaring a series known through $\Trunc{\cdot}{4}$ with
$\ordt=0$ returns a result known through $\Trunc{\cdot}{4}$, so the check is
available at every block computed and at none beyond.
\end{example}

\begin{example}[A product in the iterated field]
\label{plt:ex:mul-worked-iterated}
In $\PLit$ the coefficient convolution is finite at the inner level too.
Take the two blocks
\[
 f=\frac1{L+1}=L^{-1}-L^{-2}+L^{-3}-\cdots,
 \qquad
 g=\frac1{L-1}=L^{-1}+L^{-2}+L^{-3}+\cdots
\]
in $\K((L^{-1}))$.  For $r\ge2$, the coefficient of $L^{-r}$ in $fg$ is a sum
over the finitely many decompositions $r=i+j$ with $i,j\ge1$:
\[
 \kappa_{-r}(fg)=\sum_{i=1}^{r-1}(-1)^{i+1}\cdot1
 =\begin{cases}1,&r\ \text{even},\\0,&r\ \text{odd},\end{cases}
\]
so $fg=L^{-2}+L^{-4}+L^{-6}+\cdots=(L^2-1)^{-1}$, as it must be.  The
degree law is visible: $\degL f=\degL g=-1$, $\degL(fg)=-2$, and
$\lc_L(f)=\lc_L(g)=\lc_L(fg)=1$.  Consequently, for
$F=f\,t^{-1}+\cdots$ and $G=g+\cdots$ in $\PLit$ one gets
$\lm(FG)=t^{-1}L^{-2}$ and $\lc(FG)=1$ by
\cref{plt:thm:mul-leading-law}, with no need to compute any further block.
Observe that $\lm$ here carries a \emph{negative} power of $L$: this is
legitimate in $\PLrat$ and $\PLit$ and impossible in $\PLlau$, which is
exactly the content of \cref{plt:cor:mul-domain}.
\end{example}

\subsection{The analytic product rule}
\label{plt:sec:mul-analytic}

Everything above is formal algebra.  We now record the exact circumstances
under which the block convolution computes the asymptotic expansion of an
actual product of functions.  Nothing in this subsection asserts
convergence: an asymptotic expansion is a statement about remainders, and the
formal product theorem is used as an input, not replaced.

\begin{definition}[Polynomial--logarithmic asymptotic expansion]
\label{plt:def:mul-asymptotic}
Let $\K\subseteq\ComplexNumbers$, let $X_0>1$, and let $S=[X_0,+\infty)$ with
the real branch $L=\log X$.  A function $f\colon S\to\ComplexNumbers$ is
\emph{PL-asymptotic} to $F\in\PLlau$, written $f\approx F$, if for every
$N\in\IntegerNumbers$ there exist $M_N\in\NaturalNumbers$, $C_N>0$ and
$X_N\ge X_0$ such that
\begin{equation}\label{plt:eq:mul-asymptotic}
 \AbsoluteValue{f(X)-\bigl(\Trunc{F}{N}\bigr)(X)}
 \le C_N\,X^{-N}(\log X)^{M_N}
 \qquad(X\ge X_N),
\end{equation}
where $\bigl(\Trunc{F}{N}\bigr)(X)=\sum_{\ordt F\le n<N}p_n(\log X)X^{-n}$ is
the actual finite sum obtained by substituting the generators.  Equivalently,
$f-\Trunc{F}{N}=\BigOAt{X^{-N}(\log X)^{M_N}}{X\to+\infty,\ X\in S}$.
\end{definition}

\begin{lemma}[Basic properties of $\approx$]
\label{plt:lem:mul-asymptotic-basic}
Let $f,g\colon S\to\ComplexNumbers$ and $F,G\in\PLlau$.
\begin{enumerate}
\item If \eqref{plt:eq:mul-asymptotic} holds for arbitrarily large $N$, it
holds for every $N$; so it suffices to verify it along a cofinal set of
cutoffs.
\item $f\approx F$ and $f\approx F'$ imply $F=F'$.
\item If $f\approx F$ and $g\approx G$, then $f+g\approx F+G$ and
$\lambda f\approx\lambda F$ for $\lambda\in\K$.
\item $f\approx0$ if and only if $f=\BigOAt{X^{-N}}{X\to+\infty}$ for every
$N$.
\end{enumerate}
\end{lemma}

\begin{proof}
(1) Let $N'<N$ and suppose \eqref{plt:eq:mul-asymptotic} holds at $N$.  Then
\[
 f-\Trunc{F}{N'}=\bigl(f-\Trunc{F}{N}\bigr)
   +\sum_{N'\le n<N}p_n(\log X)X^{-n},
\]
a finite sum; the first term is $\BigOAt{X^{-N}(\log
X)^{M_N}}{X\to+\infty}\subseteq\BigOAt{X^{-N'}(\log X)^{M_N}}{X\to+\infty}$
since $N\ge N'$, and each term of the sum is
$\BigOAt{X^{-N'}(\log X)^{\degL p_n}}{X\to+\infty}$.

(2) Suppose $D=F-F'\ne0$, with $n_1=\ordt D$ and leading block $p$ of degree
$d$ and leading coefficient $c\ne0$.  Take $N=n_1+1$ and subtract the two
instances of \eqref{plt:eq:mul-asymptotic}:
\[
 \bigl(\Trunc{D}{n_1+1}\bigr)(X)=p(\log X)X^{-n_1}
 =\BigOAt{X^{-n_1-1}(\log X)^{M}}{X\to+\infty}
\]
for some $M$.  But $|p(\log X)|\sim|c|(\log X)^{d}$, so the left side
divided by $X^{-n_1-1}(\log X)^M$ is $\sim|c|X(\log X)^{d-M}\to+\infty$, a
contradiction.

(3) Immediate from
$\Trunc{F+G}{N}=\Trunc{F}{N}+\Trunc{G}{N}$ and the triangle inequality.

(4) If $f\approx0$ then $\Trunc{0}{N}=0$ and $f=\BigOAt{X^{-N}(\log
X)^{M_N}}{X\to+\infty}$ for every $N$; since $(\log
X)^{M_N}=\LittleOAt{X}{X\to+\infty}$, this gives
$f=\BigOAt{X^{-N+1}}{X\to+\infty}$ for every $N$, hence
$f=\BigOAt{X^{-N}}{X\to+\infty}$ for every $N$.  The converse is trivial.
\end{proof}

\begin{theorem}[Analytic product rule]
\label{plt:thm:mul-analytic-product}
Let $\K\subseteq\ComplexNumbers$, $S=[X_0,+\infty)$ with the real logarithm,
and let $f,g\colon S\to\ComplexNumbers$ with $f\approx F$ and $g\approx G$
for $F,G\in\PLlau$.  Then
\[
 fg\approx FG,
\]
where $FG$ is the formal Cauchy product \eqref{plt:eq:mul-cauchy}.
Consequently the set
$\mathcal A(S)=\SetOf{f:\exists F\in\PLlau,\ f\approx F}$ is a
$\K$-algebra, the map $f\mapsto F$ is a well-defined $\K$-algebra
homomorphism $\mathcal A(S)\to\PLlau$, and its kernel is the ideal of
functions that are $\BigOAt{X^{-N}}{X\to+\infty}$ for every $N$.
\end{theorem}

\begin{proof}
First dispose of the degenerate case.  Any $g$ with $g\approx G$ satisfies
$g=\BigOAt{X^{c}(\log X)^{m}}{X\to+\infty}$ for some $c\in\IntegerNumbers$
and $m\in\NaturalNumbers$: take $N=1$ in \eqref{plt:eq:mul-asymptotic}, note
that $\Trunc{G}{1}$ is a finite sum of monomials $t^nL^j$, and set
$c=\max(-\ordt G,-1)$ (with $c=-1$ if $G=0$).  Hence, if $F=0$, then
\cref{plt:lem:mul-asymptotic-basic}(4) gives
$f=\BigOAt{X^{-P}}{X\to+\infty}$ for every $P$, so
$fg=\BigOAt{X^{-P+c}(\log X)^{m}}{X\to+\infty}$ for every $P$, whence
$fg\approx0=FG$ by \cref{plt:lem:mul-asymptotic-basic}(4) again; the case
$G=0$ is symmetric.  Assume from now on $F,G\ne0$ and set $a=\ordt F$,
$b=\ordt G$.

Fix $N\in\IntegerNumbers$ and put
\[
 P=\max(N-b,\;a),
 \qquad
 Q=\max(N-a,\;b).
\]
Write $f_P=(\Trunc{F}{P})(X)$, $r=f-f_P$, $g_Q=(\Trunc{G}{Q})(X)$,
$s=g-g_Q$.  By hypothesis
$r=\BigOAt{X^{-P}(\log X)^{m_1}}{X\to+\infty}$ and
$s=\BigOAt{X^{-Q}(\log X)^{m_2}}{X\to+\infty}$ for suitable exponents.  Since
$\Trunc{F}{P}$ is a finite sum of blocks with outer exponents $\ge a$,
\[
 f_P=\BigOAt{X^{-a}(\log X)^{m_3}}{X\to+\infty},
 \qquad
 g_Q=\BigOAt{X^{-b}(\log X)^{m_4}}{X\to+\infty}
\]
with $m_3=\max_{a\le n<P}\degL p_n$ and similarly for $m_4$ (the estimates
being vacuously true if the sums are empty).  Therefore
\[
 fg-f_Pg_Q=f_Ps+rg_Q+rs
\]
is, term by term,
\[
 \BigOAt{X^{-(a+Q)}(\log X)^{m_3+m_2}}{X\to+\infty}
 +\BigOAt{X^{-(P+b)}(\log X)^{m_1+m_4}}{X\to+\infty}
 +\BigOAt{X^{-(P+Q)}(\log X)^{m_1+m_2}}{X\to+\infty}.
\]
Now $a+Q\ge a+(N-a)=N$, $P+b\ge(N-b)+b=N$, and $P+Q\ge a+(N-a)=N$; hence
\begin{equation}\label{plt:eq:mul-analytic-step1}
 fg-f_Pg_Q=\BigOAt{X^{-N}(\log X)^{M'}}{X\to+\infty}
\end{equation}
for some $M'$.

Next, $f_Pg_Q$ is the value at $X$ of the finite polynomial--logarithmic
expression $\bigl(\Trunc{F}{P}\bigr)\bigl(\Trunc{G}{Q}\bigr)\in\PLlau$,
because substitution of the generators $t=X^{-1}$, $L=\log X$ into a finite
sum of monomials is a ring homomorphism into functions on $S$.  Since
$P\ge N-b$ and $Q\ge N-a$, \cref{plt:cor:mul-truncated-product} gives
\[
 \Trunc{\bigl(\Trunc{F}{P}\bigr)\bigl(\Trunc{G}{Q}\bigr)}{N}
 =\Trunc{FG}{N}.
\]
Hence $\bigl(\Trunc{F}{P}\bigr)\bigl(\Trunc{G}{Q}\bigr)-\Trunc{FG}{N}$ is a
\emph{finite} sum of monomials $c\,t^nL^j$ with $n\ge N$, so its value at
$X$ is $\BigOAt{X^{-N}(\log X)^{M''}}{X\to+\infty}$ with $M''$ the largest
$j$ occurring.  Combining with \eqref{plt:eq:mul-analytic-step1},
\[
 fg-\bigl(\Trunc{FG}{N}\bigr)(X)
 =\BigOAt{X^{-N}(\log X)^{\max(M',M'')}}{X\to+\infty},
\]
which is \eqref{plt:eq:mul-asymptotic} for $fg$ and $FG$.  As $N$ was
arbitrary, $fg\approx FG$.

The algebra statements now follow: $\mathcal A(S)$ is closed under $+$,
scalars and $\cdot$ by \cref{plt:lem:mul-asymptotic-basic}(3) and the above;
the assignment $f\mapsto F$ is well defined by
\cref{plt:lem:mul-asymptotic-basic}(2) and is a homomorphism by the same two
statements; its kernel is described by
\cref{plt:lem:mul-asymptotic-basic}(4).
\end{proof}

\begin{remark}[Which hypotheses are actually used, and the sectorial version]
\label{plt:rmk:mul-analytic-hypotheses}
The proof uses only three things: that $f$ and $g$ have remainders bounded by
$X^{-N}$ times \emph{some} power of $\log X$; that finitely many blocks are
involved at each stage; and the formal finite-determination theorem
\cref{plt:thm:mul-finite-determination}, which is what allows the two
truncation cutoffs $P$ and $Q$ to be chosen independently of the shape of
$F$ and $G$.  No convergence and no uniformity in $L$ is used or obtained.

The same proof works verbatim on an unbounded set
$S\subseteq\SetOf{z:\AbsoluteValue{\arg z}\le\theta}$ with
$0<\theta<\pi$ and the principal logarithm, with $X^{-N}$ replaced by
$\AbsoluteValue{X}^{-N}$ and $(\log X)^{M}$ by
$(\log\AbsoluteValue{X})^{M}$; the only additional input is
$\AbsoluteValue{\log X}\le\log\AbsoluteValue{X}+\theta
\le2\log\AbsoluteValue{X}$ for $\AbsoluteValue{X}\ge\EulerE^{\theta}$, so
that a bound in $\AbsoluteValue{\log X}$ and a bound in
$\log\AbsoluteValue{X}$ differ only by a constant factor.  Uniformity claims
on a sector must be stated for a \emph{closed} subsector; the branch of
$\log$ is part of the statement and cannot be recovered from the formal
data.
\end{remark}

\begin{warningbox}
Three cautions, each corresponding to an overstatement found in the sources.
\begin{enumerate}
\item \emph{A formal identity is not an analytic theorem.}
\cref{plt:thm:mul-analytic-product} presupposes that $f$ and $g$
\emph{already have} expansions; it does not produce one.  The
homomorphism $f\mapsto F$ has a large kernel --- $\EulerE^{-X}\approx0$ ---
so a formal computation determines an actual function at best up to a flat
error.
\item \emph{The formal $\FormalBigOAt{t^N}{t}$ carries no uniformity in $L$.}
Asserting $F=G+\FormalBigOAt{t^N}{t}$ says only that the blocks agree below
outer order $N$.  In particular the constant $C_N$ and the exponent $M_N$ of
\eqref{plt:eq:mul-asymptotic} are not supplied by any formal statement and
must be produced separately.
\item \emph{The three $\BigO$-roles are distinct.}  The formal tail
$\FormalBigOAt{t^N}{t}$, the analytic estimate
$\BigOAt{X^{-N}(\log X)^{M}}{X\to+\infty}$ of
\eqref{plt:eq:mul-asymptotic}, and the complexity bound
$\BigO\bigl((N-a-b)^2d^2\bigr)$ of \cref{plt:rmk:mul-cost} are three
different assertions and none follows from the others.  A bare $\BigO$ for an
asymptotic claim, with the regime left to context, is never acceptable.
\end{enumerate}
\end{warningbox}

\begin{summarybox}
Multiplication is Cauchy convolution in $t=X^{-1}$ with exact coefficient
arithmetic in $L=\log X$.  Every output block below a cutoff is a finite
bilinear expression in finitely many input blocks, with the exact range
$a\le n\le N-b$, $b\le m\le N-a$; the resulting precision rule is
$\Trunc{FG}{\min(P+b,\,Q+a)}$ and it cannot be improved.  All four ambient
classes are closed and are integral domains, the leading monomial and leading
scalar are multiplicative, and the rank-two valuation
$v(F)=(\ordt F,-\degL p_{\ordt F})$ is additive.  The log-degree profile
obeys the max-plus bound $\dprofile{FG}\le\dprofile{F}\ast\dprofile{G}$,
with equality exactly when the maximizing pairs do not cancel; affine
profiles are stable under products, powers, and binomial powers, but not
under $t$-adic limits.  Binomial powers $(1+U)^\sigma$ are defined for every
$\sigma\in\K$ by a summable family, are finitely determined at every order,
and satisfy the exponent law.  Under printed remainder hypotheses on a stated
domain and branch, the formal product computes the asymptotic expansion of
the analytic product --- and nothing more.
\end{summarybox}
